% Target: rmp-iii-b-braid-three
% Formula: R_{j,a} replaces the upper a-j crossing in the third braid panel.
% Ink: Canonical public kernel plane patch with braid strands and an explicit
%   R_{j,a} box.
% Boundary: Eight lattice legs run past the patch and four site legs leave the
%   plane; string a enters from below and string b exits below.
% Source: paper TN-Review-main.tex line 1679; author source ImagesReview Section
%   III/ImagesReview Section III.tex lines 863-899
% Capabilities: kernel, plane-frame, braid-resolver, crossing-order
\[
  \tenkzkernel
  \tndeclare{species}{operator}{hue=source:red}
  % The author draws all four lattice lines inside one plane, so every one of
  % them carries a virtual bond; the site index is the stub that leaves the
  % plane, and the four stubs end in air.
  \begin{tenkz}[lattice={2x2}, frame=plane, physical=up,
                west=open, east=open, north=open, south=open]
    % the two-by-two lattice patch: the anyons stand on the right column
    \tn[at=(1,1), name=p]{}
    \tn[at=(1,2), name=i, role=marked, label pos=nw]{i}
    \tn[at=(2,1), name=q]{}
    \tn[at=(2,2), name=j, role=extra, label pos=sw]{j}
    % the resolver box stands east of the patch where the a-j crossing was
    \tn[at=0.9 e of midway i and j, skin=box, name=R]{\noexpand\mathcal R_{j,a}}
    % string a's former via waypoints, now junctions: three coincide with the
    % lattice-crossing marker beads (kept visible), the two bends over i are
    % new and invisible
    \tn[at=0.3 w of q, name=ja1]{}
    \tn[at=midway p and q, name=ja2]{}
    \tn[at=midway p and i, name=ja3]{}
    \tn[at=0.3 nw of i, skin=none, name=ja4]{}
    \tn[at=0.3 ne of i, skin=none, name=ja5]{}
    \tn[at=0.4 e of i, name=ja6]{}
    % string b's former via waypoints: the bends off i are new and invisible,
    % the last one is the marker bead where b pierces past j
    \tn[at=0.3 se of i, skin=none, name=jb1]{}
    \tn[at=0.7 e of midway i and j, skin=none, name=jb2]{}
    \tn[at=0.4 e of j, name=jb3]{}
    % string a winds through the patch into the resolver, lapping over the
    % lattice line that runs north out of i; every lattice crossing carries
    % its dot
    \tnwire[kind=string, species=operator, name=a1, route=arc]{open s}{ja1}
    \tnwire[kind=string, species=operator, name=a2, route=arc,
            cross={over at crossing of a2 and a1}]{ja1}{ja2}
    \tnwire[kind=string, species=operator, name=a3, route=arc,
            cross={over at crossing of a3 and a2}]{ja2}{ja3}
    \tnwire[kind=string, species=operator, name=a4, route=arc,
            cross={over at crossing of a4 and a3}]{ja3}{ja4}
    \tnwire[kind=string, species=operator, name=a5, route=arc,
            cross={over at crossing of a5 and open-north-col-2,
                   over at crossing of a5 and a4}]{ja4}{ja5}
    \tnwire[kind=string, species=operator, name=a6, route=arc,
            cross={over at crossing of a6 and a5}]{ja5}{ja6}
    \tnwire[kind=string, species=operator, name=a7, route=arc,
            cross={over at crossing of a7 and a6}]{ja6}{R}
    % string b leaves anyon i, dives behind the resolver, and exits below
    \tnwire[kind=string, species=operator, name=b1, route=arc]{i}{jb1}
    \tnwire[kind=string, species=operator, name=b2, route=arc,
            cross={over at crossing of b2 and b1}]{jb1}{jb2}
    \tnwire[kind=string, species=operator, name=b3, route=arc,
            cross={over at crossing of b3 and b2}]{jb2}{jb3}
    \tnwire[kind=string, species=operator, name=b4, route=arc,
            cross={over at crossing of b4 and b3}]{jb3}{open s}
    % anyon j hands its strand to the resolver, ending under its west edge
    \tnwire[kind=string, species=operator, name=ja,
            route=arc]{j}{0.55 e of midway i and j}
    % the marker bead where a laps over the lattice line north of i
    \tn[at=crossing of a5 and open-north-col-2]{}
  \end{tenkz}
\]
